\hypertarget{chapter-22-c17-vocabulary-types}{%
\section{CHAPTER 22: C17 VOCABULARY
TYPES}\label{chapter-22-c17-vocabulary-types}}

\hypertarget{c17-vocabulary-types}{%
\section{C++17 VOCABULARY TYPES}\label{c17-vocabulary-types}}

These types provide standard ways to represent optionality,
alternatives, and type-erased values, replacing many custom
implementations.

\hypertarget{stdstring_view}{%
\subsection{\texorpdfstring{1.
\texttt{std::string\_view}}{1. std::string\_view}}\label{stdstring_view}}

A non-owning reference to a string (or substring). \textbf{Zero-copy}
string operations.

\hypertarget{efficiency}{%
\subsubsection{1.1 Efficiency}\label{efficiency}}

\begin{lstlisting}[language={C++}]
// Bad: Copies string (potentially expensive allocation)
void print_str(std::string s) {
    std::cout << s << "\n";
}

// Good: No copy, no allocation
void print_view(std::string_view sv) {
    std::cout << sv << "\n";
}

int main() {
    const char* cstr = "Hello World";
    // print_str(cstr); // Creates std::string (allocates!)
    print_view(cstr);   // No allocation

    std::string s = "Hello World";
    print_view(s);      // Works with std::string too

    // Substrings are cheap!
    std::string_view sub = std::string_view(cstr).substr(0, 5);
    print_view(sub);
}
\end{lstlisting}

\hypertarget{caveats}{%
\subsubsection{1.2 Caveats}\label{caveats}}

\begin{itemize}
\tightlist
\item
  \textbf{Non-owning:} Ensure the underlying string outlives the view.
\item
  \textbf{Not null-terminated:} Do not pass
  \passthrough{\lstinline!.data()!} to C APIs unless you are sure.
\end{itemize}

\hypertarget{stdoptional}{%
\subsection{\texorpdfstring{2.
\texttt{std::optional}}{2. std::optional}}\label{stdoptional}}

Represents a value that may or may not be present. Replaces pointers for
nullable values or ``magic values'' (-1, "").

\begin{lstlisting}[language={C++}]
#include <optional>

std::optional<int> find_even(const std::vector<int>& v) {
    for (int x : v) {
        if (x % 2 == 0) return x;
    }
    return std::nullopt; // or {}
}

int main() {
    auto res = find_even({1, 3, 5});
    if (res) { // or res.has_value()
        std::cout << *res; // or res.value() (throws if empty)
    } else {
        std::cout << "Not found";
    }

    // Value or default
    std::cout << res.value_or(0);
}
\end{lstlisting}

\hypertarget{stdvariant}{%
\subsection{\texorpdfstring{3.
\texttt{std::variant}}{3. std::variant}}\label{stdvariant}}

A type-safe union. Can hold one of several distinct types.

\begin{lstlisting}[language={C++}]
#include <variant>

std::variant<int, float, std::string> v;

v = 10;
v = 3.14f;
v = "hello";

// Accessing
try {
    std::string s = std::get<std::string>(v);
    // int i = std::get<int>(v); // Throws std::bad_variant_access
} catch (...) {}

// std::visit (The Visitor Pattern)
std::visit([](auto&& arg) {
    std::cout << arg << "\n";
}, v);
\end{lstlisting}

\hypertarget{stdany}{%
\subsection{\texorpdfstring{4.
\texttt{std::any}}{4. std::any}}\label{stdany}}

A type-safe container for \emph{single} values of any type. (Like
\passthrough{\lstinline!void*!} but safe).

\begin{lstlisting}[language={C++}]
#include <any>

std::any a = 1;
a = std::string("hello");

try {
    std::string s = std::any_cast<std::string>(a);
} catch (const std::bad_any_cast& e) {
    std::cout << e.what();
}
\end{lstlisting}

\hypertarget{professional-notes-chapter-51-stdoptional}{%
\section{Professional Notes: Chapter 51:
std::optional}\label{professional-notes-chapter-51-stdoptional}}

Section 51.1: Using optionals to represent the absence of a value
Section 51.2: optional as return value Section 51.3: value\_or Section
51.4: Introduction Section 51.5: Using optionals to represent the
failure of a function

\hypertarget{professional-notes-chapter-56-stdvariant}{%
\section{Professional Notes: Chapter 56:
std::variant}\label{professional-notes-chapter-56-stdvariant}}

Section 56.1: Create pseudo-method pointers Section 56.2: Basic
std::variant use Section 56.3: Constructing a
\passthrough{\lstinline!std::variant!}

\hypertarget{professional-notes-chapter-51-stdoptional-1}{%
\section{Professional Notes: Chapter 51:
std::optional}\label{professional-notes-chapter-51-stdoptional-1}}

Section 51.1: Using optionals to represent the absence of a value Before
C++17, having pointers with a value of nullptr commonly represented the
absence of a value. This is a good solution for large objects that have
been dynamically allocated and are already managed by pointers. However,
this solution does not work well for small or primitive types such as
int, which are rarely ever dynamically allocated or managed by pointers.
std::optional provides a viable solution to this common problem. In this
example, struct Person is dened. It is possible for a person to have a
pet, but not necessary. Therefore, the pet member of Person is declared
with an std::optional wrapper. \#include \#include \#include struct
Animal \{ std::string name; \}; struct Person \{ std::string name;
std::optional pet; \}; int main() \{ Person person; person.name =
``John''; if (person.pet) \{ std::cout \textless\textless{} person.name
\textless\textless{} ``'s pet's name is'' \textless\textless{}
person.pet-\textgreater name \textless\textless{} std::endl; \} else \{
std::cout \textless\textless{} person.name \textless\textless{} " is
alone." \textless\textless{} std::endl; \} \} Section 51.2: optional as
return value std::optional divide(float a, float b) \{ if (b!=0.f)
return a/b; return \{\}; \} Here we return either the fraction a/b, but
if it is not dened (would be innity) we instead return the empty
optional. A more complex case: template\textless class Range, class
Pred\textgreater{} auto find\_if( Range\&\& r, Pred\&\& p ) \{ using
std::begin; using std::end; auto b = begin(r), e = end(r); auto r =
std::find\_if(b, e , p ); using iterator = decltype(r);  if (r==e)
return std::optional(); return std::optional(r); \}
template\textless class Range, class T\textgreater{} auto find(
Range\&\& r, T const\& t ) \{ return find\_if( std::forward(r),
\href{auto\&\&\%20x}{\&t}\{return x==t;\} ); \} find( some\_range, 7 )
searches the container or range some\_range for something equal to the
number 7. find\_if does it with a predicate. It returns either an empty
optional if it was not found, or an optional containing an iterator
tothe element if it was. This allows you to do: if (find( vec, 7 )) \{
// code \} or even if (auto oit = find( vec, 7 )) \{
vec.erase(\emph{oit); \} without having to mess around with begin/end
iterators and tests. Section 51.3: value\_or void print\_name(
std::ostream\& os, std::optional const\& name ) \{ std::cout ``Name
is:'' \textless\textless{} name.value\_or(``'') \textless\textless{}
`\n'; \} value\_or either returns the value stored in the optional, or
the argument if there is nothing store there. This lets you take the
maybe-null optional and give a default behavior when you actually need a
value. By doing it this way, the ``default behavior'' decision can be
pushed back to the point where it is best made and immediately needed,
instead of generating some default value deep in the guts of some
engine. Section 51.4: Introduction Optionals (also known as Maybe types)
are used to represent a type whose contents may or may not be present.
They are implemented in C++17 as the std::optional class. For example,
an object of type std::optional may contain some value of type int, or
it may contain no value. Optionals are commonly used either to represent
a value that may not exist or as a return type from a function that can
fail to return a meaningful result. Other approaches to optional There
are many other approach to solving the problem that std::optional
solves, but none of them are quite complete: using a pointer, using a
sentinel, or using a pair\textless bool, T\textgreater. Optional vs
Pointer In some cases, we can provide a pointer to an existing object
or nullptr to indicate failure. But this is limited to those cases where
objects already exist - optional, as a value type, can also be used to
return new objects without resorting to memory allocation. Optional vs
Sentinel A common idiom is to use a special value to indicate that the
value is meaningless. This may be 0 or -1 for integral types, or nullptr
for pointers. However, this reduces the space of valid values (you
cannot dierentiate between a valid 0 and a meaningless 0) and many types
do not have a natural choice for the sentinel value. Optional vs
std::pair\textless bool, T\textgreater{} Another common idiom is to
provide a pair, where one of the elements is a bool indicating whether
or not the value is meaningful. This relies upon the value type being
default-constructible in the case of error, which is not possible for
some types and possible but undesirable for others. An optional, in the
case of error, does not need to construct anything. Section 51.5: Using
optionals to represent the failure of a function Before C++17, a
function typically represented failure in one of several ways: A null
pointer was returned. e.g.~Calling a function Delegate
}App::get\_delegate() on an App instance that did not have a delegate
would return nullptr. This is a good solution for objects that have been
dynamically allocated or are large and managed by pointers, but isn't a
good solution for small objects that are typically stack-allocated and
passed by copying. A specic value of the return type was reserved to
indicate failure. e.g.~Calling a function unsigned
shortest\_path\_distance(Vertex a, Vertex b) on two vertices that are
not connected may return zero to indicate this fact. The value was
paired together with a bool to indicate is the returned value was
meaningful. e.g.~Calling a function std::pair\textless int,
bool\textgreater{} parse(const std::string \&str) with a string argument
that is not an integer would return a pair with an undened int and a
bool set to false. In this example, John is given two pets, Fluy and
Furball. The function Person::pet\_with\_name() is then called to
retrieve John's pet Whiskers. Since John does not have a pet named
Whiskers, the function fails and std::nullopt is returned instead.
\#include \#include \#include \#include struct Animal \{ std::string
name; \}; struct Person \{ std::string name; std::vector pets;
std::optional pet\_with\_name(const std::string \&name) \{ for (const
Animal \&pet : pets) \{  if (pet.name == name) \{ return pet; \} \}
return std::nullopt; \} \}; int main() \{ Person john; john.name =
``John''; Animal fluffy; fluffy.name = ``Fluffy'';
john.pets.push\_back(fluffy); Animal furball; furball.name =
``Furball''; john.pets.push\_back(furball); std::optional whiskers =
john.pet\_with\_name(``Whiskers''); if (whiskers) \{ std::cout
\textless\textless{} ``John has a pet named Whiskers.''
\textless\textless{} std::endl; \} else \{ std::cout
\textless\textless{} ``Whiskers must not belong to John.''
\textless\textless{} std::endl; \} \}

\hypertarget{professional-notes-chapter-56-stdvariant-1}{%
\section{Professional Notes: Chapter 56:
std::variant}\label{professional-notes-chapter-56-stdvariant-1}}

Section 56.1: Create pseudo-method pointers This is an advanced example.
You can use variant for light weight type erasure. template struct
pseudo\_method \{ F f; // enable C++17 class type deduction:
pseudo\_method( F\&\& fin ):f(std::move(fin)) \{\} // Koenig lookup
operator-\textgreater{}\emph{, as this is a pseudo-method it is
appropriate: template // maybe add SFINAE test that LHS is actually a
variant. friend decltype(auto) operator-\textgreater{}}( Variant\&\&
var, pseudo\_method const\& method ) \{ //
var-\textgreater{}\emph{method returns a lambda that perfect forwards a
function call, // behaving like a method pointer basically: return
\href{auto\&\&...args}{\&}-\textgreater decltype(auto) \{ // use visit
to get the type of the variant: return std::visit(
\href{auto\&\&\%20self}{\&}-\textgreater decltype(auto) \{ //
decltype(x)(x) is perfect forwarding in a lambda: return method.f(
decltype(self)(self), decltype(args)(args)\ldots{} ); \},
std::forward(var) ); \}; \} \}; this creates a type that overloads
operator-\textgreater{}} with a Variant on the left hand side. // C++17
class type deduction to find template argument of
\passthrough{\lstinline!print!} here. // a pseudo-method lambda should
take \passthrough{\lstinline!self!} as its first argument, then // the
rest of the arguments afterwards, and invoke the action: pseudo\_method
print =
\href{auto\&\&\%20self,\%20auto\&\&...args}{}-\textgreater decltype(auto)
\{ return decltype(self)(self).print( decltype(args)(args)\ldots{} );
\}; Now if we have 2 types each with a print method: struct A \{ void
print( std::ostream\& os ) const \{ os \textless\textless{} ``A''; \}
\}; struct B \{ void print( std::ostream\& os ) const \{ os
\textless\textless{} ``B''; \} \}; note that they are unrelated types.
We can: std::variant\textless A,B\textgreater{} var = A\{\};
(var-\textgreater{}\emph{print)(std::cout); and it will dispatch the
call directly to A::print(std::cout) for us. If we instead initialized
the var with B\{\}, it would dispatch to B::print(std::cout). If we
created a new type C: struct C \{\}; then:
std::variant\textless A,B,C\textgreater{} var = A\{\};
(var-\textgreater{}}print)(std::cout); will fail to compile, because
there is no C.print(std::cout) method. Extending the above would permit
free function prints to be detected and used, possibly with use of if
constexpr within the print pseudo-method. live example currently using
boost::variant in place of std::variant. Section 56.2: Basic
std::variant use This creates a variant (a tagged union) that can store
either an int or a string. std::variant\textless{} int, std::string
\textgreater{} var; We can store one of either type in it: var =
``hello''s; And we can access the contents via std::visit: // Prints
``hello\n'': visit( \href{auto\&\&\%20e}{} \{ std::cout
\textless\textless{} e \textless\textless{} `\n'; \}, var ); by passing
in a polymorphic lambda or similar function object. If we are certain we
know what type it is, we can get it: auto str = std::get(var); but this
will throw if we get it wrong. get\_if: auto* str = std::get\_if(\&var);
returns nullptr if you guess wrong. Variants guarantee no dynamic memory
allocation (other than which is allocated by their contained types).
Only one of the types in a variant is stored there, and in rare cases
(involving exceptions while assigning and no safe way to back out) the
variant can become empty. Variants let you store multiple value types in
one variable safely and eciently. They are basically smart, type-safe
unions. Section 56.3: Constructing a
\passthrough{\lstinline!std::variant!} This does not cover allocators.
struct A \{\}; struct B \{ B()=default; B(B const\&)=default;
B(int)\{\}; \}; struct C \{ C()=delete; C(int) \{\}; C(C
const\&)=default; \}; struct D \{ D( std::initializer\_list ) \{\}; D(D
const\&)=default; D()=default; \};
std::variant\textless A,B\textgreater{} var\_ab0; // contains a A()
std::variant\textless A,B\textgreater{} var\_ab1 = 7; // contains a B(7)
std::variant\textless A,B\textgreater{} var\_ab2 = var\_ab1; // contains
a B(7) std::variant\textless A,B,C\textgreater{} var\_abc0\{
std::in\_place\_type, 7 \}; // contains a C(7) std::variant var\_c0; //
illegal, no default ctor for C std::variant\textless A,D\textgreater{}
var\_ad0( std::in\_place\_type, \{1,3,3,4\} ); // contains D\{1,3,3,4\}
std::variant\textless A,D\textgreater{} var\_ad1(
std::in\_place\_index\textless0\textgreater{} ); // contains A\{\}
std::variant\textless A,D\textgreater{} var\_ad2(
std::in\_place\_index\textless1\textgreater, \{1,3,3,4\} ); // contains
D\{1,3,3,4\}
